\section{Related Literature and Contribution}
\label{sec:lit}

This paper relates to three strands of research: the literature on local interaction, agglomeration, and regional coordination; the literature on proximity, lock-in, and system renewal; and the welfare analysis of place-based policy with particular attention to external linkages and regional policy design. The paper's contribution is to connect these strands in a two-region dynamic welfare framework in which a public hub raises local coordination within a target thin region, while a complementary pipeline policy creates cross-regional renewal through outside linkages.

\subsection{Local Interaction, Agglomeration, and Regional Coordination}

A first relevant literature explains why spatial concentration can raise productivity. In urban and regional economics, agglomeration economies are commonly rationalized through mechanisms such as matching, learning, and sharing \citep{DurantonPuga2004,RosenthalStrange2004}. In this view, spatial proximity improves the efficiency with which actors find collaborators, exchange information, and combine complementary ideas or capabilities. This logic provides the most direct theoretical motivation for a public co-creation hub: a hub can be understood as a policy device that raises the effective rate at which potentially valuable local interactions are realized within a region.

Within this broader perspective, face-to-face interaction has been emphasized as a distinct source of value. \citet{StorperVenables2004} argue that face-to-face contact facilitates trust formation, complex communication, and the transfer of tacit knowledge. \citet{JaffeTrajtenbergHenderson1993} show that knowledge spillovers are geographically localized, while \citet{RocheOettlCatalini2024} provide more direct evidence that proximate co-working environments can generate knowledge spillovers through local social interaction. These contributions support the paper's premise that reducing local interaction frictions may create social value even in relatively small regional ecosystems.

The paper also speaks to the emerging literature on co-working spaces, enterprise hubs, and collaborative spaces outside major metropolitan cores. \citet{VoglAkhavan2022} review the evidence on co-working spaces in peripheral and rural areas and conclude that their effects are heterogeneous and context-dependent. \citet{MerrellPhillipsonGortonCowie2022} analyze enterprise hubs as a mechanism for rural local economic development, while \citet{MerrellRoweCowieGkartzios2021} interpret rural enterprise hubs as micro-clusters that may support creativity-led development. This literature is close to the policy object studied here, but it remains largely empirical and descriptive. It documents possible benefits of hubs and the conditions under which they succeed, but it does not offer a welfare criterion for determining when a publicly supported hub is efficient, when it merely generates visible short-run activity, or when it should be complemented by external-linkage policy.

Relative to these literatures, the present paper does not simply restate that local interaction is valuable. Instead, it embeds a hub in a regional welfare framework in which local coordination is beneficial but not automatically sufficient for dynamic efficiency. In that sense, the paper takes the interaction-based logic of agglomeration seriously, but shifts the question from whether local interaction matters to when a public device that organizes local interaction is socially justified.

\subsection{Proximity, Lock-In, and System Renewal}

A second relevant literature emphasizes that proximity has both benefits and costs. \citet{Boschma2005} argues that proximity can support innovation, but that too much proximity can generate lock-in by narrowing the scope for new combinations. \citet{TorreRallet2005} similarly distinguish different dimensions of proximity and stress that spatial closeness is neither universally beneficial nor sufficient for sustained innovation. This literature is central for the present paper because the model's dynamic redundancy mechanism is precisely a formalization of the idea that repeated local interaction may become less novel over time.

The paper is also closely related to the regional innovation systems literature. \citet{CookeUrangaEtxebarria1997} and \citet{AsheimLawtonSmithOughton2011} stress that regional innovation depends on the institutional and organizational structure of local systems rather than on scale alone. \citet{TodtlingTrippl2005} argue that innovation policy should be differentiated across places, while \citet{CoenenAsheimBuggeHerstad2017} emphasize evolutionary and system-level perspectives in regional policy. Work on path renewal and path creation further shows that even regions with established specializations may struggle to generate new trajectories without appropriate renewal mechanisms \citep{ChaminadeBellandiPlecheroSantini2019}.

What matters for the present paper is that this literature does more than warn against excessive proximity. It also points toward the importance of renewal mechanisms that reconnect regions to outside sources of knowledge, collaborators, and opportunities. \citet{BatheltMalmbergMaskell2004} capture this especially well through the distinction between local buzz and global pipelines. That distinction is highly relevant here. In the model developed below, the public hub plays the role of a local coordination device inside the target region, while the pipeline policy captures a distinct margin of adjustment: the formation of extra-regional links that replenish novelty and reduce the danger of local lock-in.

This is the key conceptual bridge between the regional innovation systems literature and the present model. The model does not treat renewal as a vague narrative possibility. Rather, it treats renewal as a policy-relevant margin that can be compared directly with local coordination in welfare terms. Put differently, the paper links the literature on proximity and lock-in to an explicit criterion for deciding when local interaction should be supplemented by cross-regional renewal.

\subsection{Place-Based Policy, External Linkages, and This Paper's Contribution}

A third relevant literature studies the welfare properties of geographically targeted policy. \citet{KlineMoretti2014ARE} provide a general welfare-theoretic treatment of local economic development programs, emphasizing that place-based interventions may correct localized distortions but need not be efficient in general. \citet{KlineMoretti2014QJE} show in a major historical application that geographically targeted interventions can produce long-run local gains, while also making clear that the welfare case for such interventions depends on the underlying externalities and margins of adjustment. More broadly, \citet{NeumarkSimpson2015} survey place-based policies and stress that their effects are heterogeneous and context-dependent.

The present paper belongs to this place-based policy tradition, but studies a policy instrument that differs from the tax incentives, subsidies, and infrastructure programs more commonly analyzed in that literature in that the object here is a public hub whose purpose is to reduce local interaction frictions in a thin target region. More importantly, the paper argues that the relevant policy object is often not the hub alone, but a \emph{policy bundle} that combines local coordination with cross-regional linkage formation. This is where the two-region structure matters. Region $A$ is the target thin region in which the hub raises local matching; region $B$ is an outside pool from which additional collaborations, ideas, and partnerships may arrive. The pipeline policy is therefore not an abstract add-on, but a reduced-form representation of extra-regional linkage formation, external brokerage, and system-renewal mechanisms.

This perspective also sharpens the rationale for public intervention. A short-run private or local evaluation rule may focus on visible activity generated by the hub inside region $A$, such as meetings, participation, or local matches. A social planner, by contrast, evaluates the hub more broadly, incorporating both the dynamic consequences of repeated local ties and the cross-regional spillovers generated by outside linkages. The paper's contribution is to formalize this divergence without requiring a full decentralized equilibrium model. In the model, a hub may look attractive under a short-run local criterion yet be socially undesirable once dynamic redundancy is taken into account. Conversely, a hub may fail a short-run private test and still be socially desirable once paired with cross-regional pipeline policy.

Against this background, the paper makes three contributions. First, it introduces an explicitly regional distinction between \emph{local coordination} and \emph{cross-regional renewal}. The hub raises local matching within the target region, whereas the pipeline policy creates or intensifies extra-regional links with an outside pool. Second, it provides a welfare criterion for evaluating hub policy as a \emph{policy bundle}, thereby connecting local buzz, lock-in, and external renewal within a single framework. Third, it generates empirically testable predictions on repeated ties, outside linkages, and participant composition. In that sense, the paper contributes not only to the theory of place-based policy, but also to the empirical agenda on how hubs and related coordination devices should be evaluated in thin regional ecosystems.

The paper therefore contributes to the agglomeration literature by formalizing a public interaction device, to the proximity and regional innovation systems literatures by connecting lock-in and renewal to explicit welfare comparisons, and to the place-based policy literature by showing why the efficient object of evaluation may be a bundled regional policy rather than a stand-alone local project.
